\paragraph{Cellular support of the central transport.}
The preceding recurrence also determines the finite geometry of the
trajectory. We consider the positions immediately before each
of the \(108\) steps; the coordinates belong to
\(D_9=\{1,\ldots,9\}\), with cyclic transport modulo nine.

\begin{proposition}[Three cyclic diagonals of the central section]
\label{prop:el-soporte-27}
The cellular projection of the trajectory starting at
\((r,c,h,v)=(5,5,1,1)\) has exact support
\begin{equation}
 \begin{aligned}
 \operatorname{Supp}_{\mathrm{cell}}(\Gamma_e)&=\{(r,c)\in D_9^2:
                 c-r\equiv0,1,-1\pmod9\},
 \\
 |\operatorname{Supp}_{\mathrm{cell}}(\Gamma_e)|&=27.
 \end{aligned}
 \label{eq:el-soporte-27}
\end{equation}
The conjugate initial orientation \((h,v)=(-1,-1)\), with the same
reversal schedule, produces the same support.
\end{proposition}
\begin{proof}
Within a block of \(27\) steps, write \(h=v=s\), where
\(s\in\{-1,1\}\). If a three-step block starts at \((a,a)\),
the additive reading precedes the displacement to \((a,a+s)\);
the multiplicative reading precedes the displacement to
\((a+s,a+s)\); the threshold preserves this last position. The three
pre-step states therefore belong to the diagonals
\(c-r=0\) and \(c-r=s\) of the positional torus. The nine
consecutive blocks traverse all nine possibilities for \(a\), since
translation by \(s\) has order nine. The subsequent reversal
traverses the adjacent diagonal of opposite sign. The diagonals
with differences \(0,1,-1\) are disjoint, each contains nine
cells, and all are visited. Reversing the initial sign
reverses the order in which the two adjacent diagonals are visited.
\end{proof}

The fixed point \((5,5)\) characterizes the section of origin; the set
\(\operatorname{Supp}_{\mathrm{cell}}(\Gamma_e)\) characterizes its cellular trajectory. The result
determines a discrete morphology on the APP domain. The spatial
realization and the orientation gluing additionally receive their
corresponding transport data. The persistence of the section and the
variation of its observable position are thus compatible properties.

\paragraph{Two central histories in the integral chart.}
The chart \(\mathcal L_{12}\) of
Section~\ref{subsec:k-registro-integral} preserves the events and their
order. We apply it to the two central orientations. For each
step \(t\), let \(d_t=\rho_9(u_t)\), where \(u_t\) is the additive,
multiplicative, or threshold integer read by the recurrence.
In the window \(E_{i+1}=\{9i+1,\ldots,9i+9\}\), define
\[
 n_i=\#\{t\in E_{i+1}:d_t\in\{2,4,6,8\}\},\qquad
 e_i=\Sigma_{12}(i+1)n_i,\qquad0\le i<12.
\]
The sign \(\Sigma_{12}\) is the sectorial orientation of the chart.
For the \({++}\) trajectory, it coincides with the transported
horizontal sign; for the \({--}\) trajectory, it coincides with its opposite.
Equivalently, for both trajectories it is
\(\Sigma_{12}(i+1)=h^{\mathrm{pre}}_{9i+1}/h_0\), where
\(h^{\mathrm{pre}}_t\) is the orientation before the step is read.
This normalization separately preserves the initial orientation and
the relative orientation. The symbols \(e_i\) denote integer
events, whereas \(e\) without a subscript denotes Euler's number.

Evaluating the two recurrences by windows gives
\[
\begin{array}{c|rrrrrrrrrrrr}
i&0&1&2&3&4&5&6&7&8&9&10&11\\\hline
n_i^{++}&2&2&5&4&3&2&2&2&5&4&3&2\\
n_i^{--}&4&3&2&2&2&5&4&3&2&2&2&5
\end{array}
\]
and determines the histories
\begin{equation}
 e_{++}=(2,2,5,-4,-3,-2)^2,\qquad
 e_{--}=(4,3,2,-2,-2,-5)^2.
 \label{eq:el-historias-centrales}
\end{equation}
The exponent denotes concatenation. Both histories have zero sum,
the same sequence of sectorial signs, and the same absolute-value
histogram: six twos, two threes, two fours, and two fives.

\begin{proposition}[Reversible separation of the two histories]
\label{prop:el-memorias-distintas}
In the coordinates \((s_C,M,a)=\mathcal L_{12}(e)\), both histories
satisfy \(s_C=0\) and \(a=0\), whereas
\begin{equation}
 M_{++}=(8,8,14,-4,-2),\qquad
 M_{--}=(18,16,14,6,6).
 \label{eq:el-memorias-distintas}
\end{equation}
The chart reconstructs each history uniquely.
\end{proposition}
\begin{proof}
The two sextuples repeat, so that
\(p_i=e_i+e_{i+6}=2e_i\) and \(a_i=e_i-e_{i+6}=0\).
Subtracting \(p_5\) from the first five \(p_i\) gives the margins
in \eqref{eq:el-memorias-distintas}. Their sums are \(24\) and
\(60\), respectively; the inverse formula in
Lemma~\ref{lem:k-carta-integral} gives \(p_5=-4\) and \(p_5=-10\).
Then \(p_i=p_5+M_i\) and
\(e_i=e_{i+6}=p_i/2\) recover the two sextuples exactly.
Divisibility by six and the six parity congruences characterizing
the integral image of the reader are satisfied.
\end{proof}

The equality of counts and signs preserves a common projection; the
margins preserve the difference between the histories. The opposite
recurrence emits \(\gamma_{--}=(b^3a^3)^2\): after \(27\) steps,
the \({++}\) route reaches the centre again with the initial
orientation of \({--}\), and the three-step schedule retains its phase.
Thus \(\gamma_{--,t}=\gamma_{e,t+27}\), with cyclic indices.
The counts \(D_k\) are invariant under this shift; the window reader
receives the same sectorial word for both histories. Both return
readers therefore coincide for these records distinguished by
the integral chart. Mass evaluation and
event reconstruction have precise informational scopes
within the same transport.
